\chapter{Applications}

\todo[inline]{lots of blabla}

\section{Ax-Grothendieck}

On présente dans un premier temps le cas particulier de $\bC^n$ et sa preuve
détaillée avant de passer au cas général d'une variété algébrique sur un corps
algébriquement clos dont l'idée est la même.

On va avoir besoin du lemme suivant: 

\begin{lemm} \label{lem-phrase-corps}
Pour tout couple $(n,d)$ d'entiers positifs, il existe une phrase
$\varphi_{n,d}$ dans le language des anneaux qui exprime que dans un corps $K$,
toute famille $(f_1, \ldots, f_n)$ de polynômes en $n$ variables de degré
inférieur ou égal à $d$, si l'application associée est injective, alors elle est surjective.
\end{lemm}

\begin{theo}[Ax-Grothendieck]\label{th-Ax-Grothendieck}
Soit $f : \bC^n \rightarrow \bC^n$ une application polynomiale. Si $f$ est
injective, alors elle est bijective. 
\end{theo}

\begin{proof}
D'après le \cref{lem-phrase-corps} et par principe de Lefschetz (\cref{coro-Lefschetz}), il suffit de
montrer le résultat
pour $\Fpalg$ pour tout $p \in \bP$. Par le \cref{th-structure-corps-fini}, toute partie finie de $\Fpalg$ est contenue dans
un sous-corps fini de $\Fpalg$. Soit $f = (f_1, \ldots, f_n) : \Fpalg^n
\rightarrow \Fpalg^n$ ,  une application polynomiale injective. Supposons qu'il
existe $y \in \Fpalg^n$ qui n'est pas dans l'image de $f$. Soit $K$ le sous-corps
fini de $\Fpalg$ qui contient toutes les composantes de $y$ et tous les
coefficients des $f_i$, alors $f$ induit une application $K^n \rightarrow K^n$
qui est injective et non surjective, ce qui est absurde car $K$ est fini. 
\end{proof}


\begin{lemm} \label{lem-phrase_variete}
Pour tout quadruplet d'entiers $(n, m, d, \delta)$, il existe une phrase $\varphi_{n,m,d,\delta}$ dans le langage des anneaux qui exprime la propriété suivante :
Pour tout corps $K$, pour toute variété algébrique $V \subset K^n$ définie par $m$ équations de degré au plus $\delta$, et pour toute application polynomiale $f: V \to V$ définie par des polynômes de degré au plus $d$, si $f$ est injective alors elle est surjective.
\end{lemm}

\begin{theo}[Ax-Grothendieck]\label{th-Ax-Grothendieck-variete}
Soit $K$ un corps algébriquement clos, $V$ une variété algébrique sur $K$ et $f
: V \rightarrow V$ une application polynomiale. Si $f$ est
injective, alors elle est bijective. 
\end{theo}

\begin{proof}
La preuve suit exactement le même schéma que pour le cas affine $V=\bC^n$ (\cref{th-Ax-Grothendieck}).
Par le lemme et par principe de Lefschetz, il suffit de démontrer le résultat pour $K=\Fpalg$.

La seule adaptation nécessaire concerne la construction du sous-corps fini. Si
l'on suppose qu'il existe $y \in V$ hors de l'image de $f$,
on considère le sous-corps fini $k \subset \Fpalg$ engendré par: les
coefficients de l'application $f$ et les coordonnées de $y$ (comme
précédemment) ET les coefficients des polynômes définissant la variété $V$.

Alors l'ensemble des points rationnels $V(k) = V \cap k^n$ est fini (car inclus dans $k^n$).
Par stabilité, $f$ induit une application $f: V(k) \to V(k)$ injective et non
surjective, ce qui est exclu.
\end{proof}




\section{Nullstellensatz}

